\chapter{系统实现与测试验证}

\section{工程实现的收束方式}

前五章已经把软件边界、总体结构、通信处理和测试算法分别说明清楚。工程实现阶段关注这些设计能否在同一个程序中按顺序运行。实际测试时，操作人员先选择串口和协议，再填写测试参数；后端随后发送控制帧、接收反馈帧，并把计算后的结果交给界面和保存模块。第六章讨论这条已经落地的执行链路。

当前实现沿用前文确定的测试类型，并把各模块接入同一运行路径。渲染进程只处理串口选择、参数填写、测试启动和结果查看；预加载脚本限定界面侧可调用的本地能力；主进程保存串口对象、测试核心和文件操作入口。通信模块负责构造协议帧、写入串口、解析反馈并维护最近样本，测试核心再使用这类样本完成通信测试、预测试、时域测试和频域测试。结果保存模块从共享结果状态读取数据，用于当前结果留档和历史结果恢复。

这种实现方式与第三章的数据流划分保持一致。高频串口反馈不直接进入界面组件，测试任务也不直接操作文件系统。控制、采集、显示和结果管理之间保持清楚边界，是桌面测控软件和工业测试系统实现中的基本要求\cite{zhang2017qtflowmeter,liu2022industrialsoftware}。本章只讨论当前已经验证的工程能力，长期运行、端到端时延和自动报告生成仍放在后续工作中说明。

\section{关键实现环节}

测试入口的实现先处理参数合法性。通信测试检查测试次数和发送间隔，预测试检查采样时长、置信水平和容许误差，时域与频域测试继续检查噪声、调节时间、平均速度、扫频范围和激励幅值等参数。参数通过后，测试核心创建任务标识并开始运行。界面侧只接收阶段、进度、统计量、完成结果和错误信息。

串口写入是另一个集中处理的位置。手动控制、自动控制、使能检测和正式测试都会产生控制量，分散写入端口会削弱控制顺序的可追踪性。当前实现把控制帧生成、发送队列、反馈解析和错误上报放在通信链路内。协议一和协议二的差别保留在帧构造与解析函数中；时域、频域等测试核心只读取角度形式的反馈样本，不再处理同步字、帧长或校验字段。

持续运行的任务由后台侧承担。自动控制帧按设定节拍生成，离线解析和渲染数据整理也不放在界面组件内执行。调度层需要记录当前处于手动模式还是自动模式，并在任务取消、串口断开或测试结束时清理相应状态。第三章提出的高频反馈流与低频交互流分离，在实现中主要落在这些状态切换和消息转发上。

结果管理的处理相对直接。通信测试保存收发统计和反馈间隔，预测试保存噪声标准差、保守调节时间和平均速度，时域和频域测试保存曲线数据、分组统计和指标摘要。保存文件重新加载后，结果被写回共享状态，图表和结果面板再按当前数据刷新。图片导出口虽然已经存在，但当前验证范围只覆盖结构化结果保存与历史加载。

\section{分层验证组织}

工程验证分成两类。第一类在脱离真实硬件的条件下运行，主要检查协议解析、统计计算、测试路径生成、频域幅相估计、控制模式切换和结果回放。第二类接入真实串口设备和舵机控制器，重点观察控制帧是否写入串口、反馈帧是否被解析、使能判定是否与实际接线状态一致。工业软件验证通常需要同时关注局部逻辑和运行环境\cite{xu2021industry4test}，本系统也按这个思路安排检查。

自动化测试的价值在于可重复。协议字段、校验逻辑、采样规模计算和结果对象结构可以反复运行，不依赖现场设备状态。实机检查的价值则在于覆盖真实链路。串口线缆、控制器反馈、发送节拍和机械状态会影响最终结果，这些因素无法通过单元测试完全替代。第六章的验证结论同时引用自动化测试和实机检查，但不把它扩展为长期稳定性结论。

\section{实机检查与结果分析}

现有实机检查只对应当前联调环境。综合检查报告中通过的项目有三项：手动链路、高频自动链路和使能检测。它们不能代表所有测试工况，却覆盖了后续测试最先依赖的三个位置，即基础收发、连续发送和执行对象状态判定。

手动链路检查时，程序持续写入控制帧，并同步读取控制器反馈。检查记录中发送失败计数为零，反馈帧能够解析为有效数据。由此可以确认，在当前协议配置下，基础串口收发路径已经跑通。

高频自动链路检查主要看后台任务生成的控制帧能否进入真实串口发送流程。报告中的转发数量、实际发送数量和失败统计之间没有矛盾，说明自动控制任务与串口发送队列之间的衔接有效。频域测试后续依赖连续激励信号，这一结果为连续控制发送提供了实机依据。

使能检测检查以当前硬件接线状态作为对照。报告中的四路状态与实际接线一致，连续轮次中未出现检测失败。该结果说明基线采样、小扰动激励和相关判别在这次联调条件下能够给出有效状态判断。

结合前端功能运行结果，当前软件已经能完成从参数输入到结果保存的基本过程。通信测试可以产生链路统计，预测试可以生成时域和频域测试所需参数，时域和频域测试可以给出图表与指标，结构化文件也可以重新加载到显示区域。该结论只适用于当前验证范围。24 h 连续运行、完整端到端时延、图片导出和长期批量归档尚未完成专门验证。

\section{本章小结}

本章说明了软件实现后各模块的运行关系。界面只发起受控请求，主进程集中管理串口、测试核心和文件操作，通信模块屏蔽协议差异，测试核心处理反馈样本并输出结果。第四章和第五章中的通信设计、预测试参数、时域指标和频域指标由同一程序承接。

验证结果来自自动化测试和当前实机检查。自动化测试覆盖协议、测试核心和控制状态切换等程序逻辑；实机检查覆盖手动收发、高频自动发送和使能检测。基于这些材料，本论文可以给出阶段性的工程验证结论。长期稳定性、完整端到端时延和图片导出流程仍需要后续补充。
